\subsection*{System}
Episode Alerts is a mobile companion app for discovering and tracking television shows.

\subsection*{Assumptions}
\begin{itemize}
	\item TMDB is available when fresh show data is requested, or cached data is already stored on the device.
	\item The user has granted notification permissions if episode reminders are enabled.
	\item Firebase is configured only when optional cloud synchronization is being used.
	\item Local device storage is available for watchlist data, preferences, and cached content.
\end{itemize}

\subsection*{The Main Use Cases Are}
\begin{itemize}
	\item Discover TV Shows
	\item Search for TV Shows
	\item View Show Details
	\item Manage Watchlist
	\item Configure Preferences and Notifications
	\item Sync Data to Cloud
\end{itemize}

\subsection*{1. Use Case 1: Discover TV Shows}
\noindent\textbf{1.1 Primary Actor:} User

\noindent\textbf{1.2 Precondition:} The app is open and the Home screen is available.

\noindent\textbf{1.3 Scenario:} The user browses the Home screen to explore featured and trending content.

\subsubsection*{1.3.1 Main Success Scenario}
\begin{enumerate}[label=\arabic*.]
	\item The user opens the Home screen.
	\item The system loads featured, airing today, popular, and top-rated shows from TMDB or from cached data.
	\item The system displays the show sections in a scrollable layout.
	\item The user taps a show card.
	\item The system opens the show details screen for the selected show.
\end{enumerate}

\subsubsection*{1.3.2 Exception Scenario}
\begin{itemize}
	\item \textbf{2a.} If step 2 fails because the network is unavailable, the system shows stale cached data and a stale-data indicator.
	\item \textbf{4a.} If step 4 cannot be completed because the selected item is unavailable, the system stays on the Home screen.
\end{itemize}

\subsection*{2. Use Case 2: Search for TV Shows}
\noindent\textbf{2.1 Primary Actor:} User

\noindent\textbf{2.2 Precondition:} The Search screen is available and the user can enter text.

\noindent\textbf{2.3 Scenario:} The user enters a query to find a specific TV show.

\subsubsection*{2.3.1 Main Success Scenario}
\begin{enumerate}[label=\arabic*.]
	\item The user opens the Search screen.
	\item The user types a show title or keyword.
	\item The system waits for a valid query and sends the request to TMDB.
	\item The system displays the matching search results.
	\item The user selects one of the results.
	\item The system opens the show details screen for that result.
\end{enumerate}

\subsubsection*{2.3.2 Exception Scenario}
\begin{itemize}
	\item \textbf{3a.} If step 3 fails because the request is too short, the system clears the results and waits for a valid search term.
	\item \textbf{3b.} If the network request fails, the system shows an error message and allows the user to retry.
	\item \textbf{4a.} If no results are returned in step 4, the system shows an empty state instead of a result list.
\end{itemize}

\subsection*{3. Use Case 3: View Show Details}
\noindent\textbf{3.1 Primary Actor:} User

\noindent\textbf{3.2 Precondition:} The user has selected a show from Home, Search, or Watchlist.

\noindent\textbf{3.3 Scenario:} The user opens a show to review detailed information and explore seasons or episodes.

\subsubsection*{3.3.1 Main Success Scenario}
\begin{enumerate}[label=\arabic*.]
	\item The user opens a show details screen.
	\item The system loads show metadata such as overview, rating, status, genres, and network information.
	\item The system displays next episode information, last episode information, and available seasons.
	\item The user opens a season to view its episode list.
	\item The user returns to the show details screen or continues to another related action, such as adding the show to the watchlist.
\end{enumerate}

\subsubsection*{3.3.2 Exception Scenario}
\begin{itemize}
	\item \textbf{2a.} If step 2 fails, the system falls back to cached details when available.
	\item \textbf{3a.} If a show does not have a next episode yet, the countdown area is hidden instead of showing invalid information.
	\item \textbf{4a.} If season data cannot be loaded, the system shows an error or empty state for that section only.
\end{itemize}

\subsection*{4. Use Case 4: Manage Watchlist}
\noindent\textbf{4.1 Primary Actor:} User

\noindent\textbf{4.2 Precondition:} The user is viewing a show card, show details screen, or the Watchlist screen.

\noindent\textbf{4.3 Scenario:} The user adds, removes, or organizes shows in the watchlist.

\subsubsection*{4.3.1 Main Success Scenario}
\begin{enumerate}[label=\arabic*.]
	\item The user taps the watchlist action for a show.
	\item The system saves the show locally in the watchlist store.
	\item The system updates the watchlist screen and related counts immediately.
	\item The system schedules release notifications when notifications are enabled.
	\item The user later removes the show or clears the watchlist if needed.
\end{enumerate}

\subsubsection*{4.3.2 Exception Scenario}
\begin{itemize}
	\item \textbf{1a.} If step 1 targets a show that is already in the watchlist, the system leaves the list unchanged and returns a duplicate-state response.
	\item \textbf{2a.} If step 2 fails because storage is unavailable, the system shows an error and keeps the previous watchlist state.
	\item \textbf{5a.} If the user clears the watchlist and later chooses undo, the system restores the last saved snapshot when available.
\end{itemize}

\subsection*{5. Use Case 5: Configure Preferences and Notifications}
\noindent\textbf{5.1 Primary Actor:} User

\noindent\textbf{5.2 Precondition:} The Settings screen is available.

\noindent\textbf{5.3 Scenario:} The user changes app preferences such as theme, analytics, and notification behavior.

\subsubsection*{5.3.1 Main Success Scenario}
\begin{enumerate}[label=\arabic*.]
	\item The user opens the Settings screen.
	\item The user changes the theme, analytics preference, or notification setting.
	\item The system saves the preference locally and updates the app behavior.
	\item If notifications are enabled, the system requests permission and schedules upcoming episode reminders.
	\item The updated setting remains active the next time the app is opened.
\end{enumerate}

\subsubsection*{5.3.2 Exception Scenario}
\begin{itemize}
	\item \textbf{2a.} If a preference update cannot be saved, the system shows an error and retains the previous value.
	\item \textbf{3a.} If analytics are disabled, the system stops recording usage events immediately.
	\item \textbf{4a.} If notification permission is denied in step 4, the system keeps notifications disabled and does not schedule reminders.
\end{itemize}

\subsection*{6. Use Case 6: Sync Data to Cloud}
\noindent\textbf{6.1 Primary Actor:} Signed-in User

\noindent\textbf{6.2 Precondition:} Firebase is configured and the user is signed in.

\noindent\textbf{6.3 Scenario:} The system synchronizes watchlist data and user preferences with the cloud.

\subsubsection*{6.3.1 Main Success Scenario}
\begin{enumerate}[label=\arabic*.]
	\item The user signs in or triggers a synchronization event.
	\item The system builds a payload from the watchlist, watch progress, watch history, and selected preferences.
	\item The system uploads the payload to Firestore.
	\item The system stores the last successful sync time.
	\item The system repeats synchronization automatically when local data changes or when the app resumes.
\end{enumerate}

\subsubsection*{6.3.2 Exception Scenario}
\begin{itemize}
	\item \textbf{1a.} If step 1 occurs while the device is offline, the system skips the upload and retries later.
	\item \textbf{3a.} If Firebase is not configured or the user is not signed in, the system keeps all data local.
	\item \textbf{5a.} If the cloud request fails, the system preserves local data and logs the sync failure for later retry.
\end{itemize}